% Standalone LaTeX source. Intentionally not compiled for this delivery.
% The inherited docs/report.tex is preserved unchanged.
\documentclass[11pt,a4paper]{article}
\usepackage[T1]{fontenc}
\usepackage[utf8]{inputenc}
\usepackage[margin=25mm]{geometry}
\usepackage{booktabs}
\usepackage{longtable}
\usepackage{listings}
\usepackage{xcolor}
\usepackage{hyperref}
\usepackage{amsmath}
\usepackage{enumitem}
\hypersetup{colorlinks=true,linkcolor=blue!45!black,urlcolor=blue!60!black,
 pdftitle={bob: An FPGA Built Inside an FPGA}}
\definecolor{codebg}{gray}{0.96}
\lstset{basicstyle=\ttfamily\scriptsize,backgroundcolor=\color{codebg},
 breaklines=true,showstringspaces=false,columns=fullflexible,keepspaces=true,
 frame=single,tabsize=4}
\setlength{\emergencystretch}{3em}
\setlist{itemsep=3pt,topsep=5pt}
\newcommand{\code}[1]{\texttt{\detokenize{#1}}}
\newcommand{\src}[1]{\par\smallskip\noindent\textbf{Source: }\nolinkurl{/Users/sk/work/bob/bob_full_v1/#1}\par\smallskip}
\title{\textbf{bob: An FPGA Built Inside an FPGA}\\[3mm]
\large Architecture, RTL, Configuration Bits, Routing, Toolchain,\\
Verification and Implementation Status on the PYNQ-Z2}
\author{Technical project report\\\texttt{bob\_full\_v1}}
\date{Source snapshot reviewed: September 2026}
\begin{document}
\maketitle
\begin{abstract}
Bob is a configurable FPGA architecture implemented as an RTL overlay in the
programmable logic of a PYNQ-Z2 board. The host device is an AMD/Xilinx Zynq
XC7Z020; the guest device is a smaller FPGA designed in Verilog and SystemVerilog.
The active profile contains sixteen configurable logic blocks, two true-dual-port
memories, two DSP-style arithmetic slices, twenty logical pads, and programmable
routing generated from a VPR routing-resource graph. A 4216-bit scan chain
configures the guest through a soft JTAG TAP driven by a Raspberry Pi Pico.

This report follows the implementation from Boolean functions to configuration
fields, and from application Verilog to simulation and hardware loading. It
separates implemented behaviour, recorded results, and proposed improvements.
It does not claim that a new build or hardware acceptance test was performed.
\end{abstract}
\tableofcontents
\newpage

\section{Scope and evidence}
\label{sec:scope}
This report describes the active sixteen-CLB M7 profile and M8 synthesis flow.
The historical report is preserved at
\nolinkurl{/Users/sk/work/bob/bob_full_v1/docs/report.tex}; its 2896-bit baseline
chain is not the current 4216-bit chain. Some historical plan sections describe
the frozen forty-eight-CLB profile. Current generated parameters and the device
schema determine the numerical descriptions used here.

\begin{description}
\item[Implemented] Present in inspected source; not necessarily accepted on the
current hardware build.
\item[Recorded simulation result] Documented in the plan or existing logs; not
rerun during report preparation.
\item[Recorded hardware result] Present in the hardware-test history or copied
Vivado reports, with milestone identity where available.
\item[Proposed] An engineering recommendation, not a silently applied change.
\end{description}
The inspected local archive contains M0--M4 Vivado reports. The hardware log
contains a passing M6 regression, including BRAM and DSP checks, but M5/M6
utilisation reports have not been copied back. M7 attempts returning
\code{0xFFFFFFFF} stop at identification, not at a functional fabric test.
A user-reported M7 WNS near $-1102$ ns awaits inspection of its actual timing
report; it is not independently verified local measurement in this document.
\src{PLAN.md}
\src{hw/src/generated/bob_params.vh}
\src{docs/hwtest/results.log}

\section{What is being built?}
\label{sec:concept}
An FPGA implements circuits using programmable truth tables, storage elements,
arithmetic blocks, memories and selectable connections. Building bob means
implementing those structures and their configuration tools, not manufacturing
new silicon. This is a \emph{soft FPGA}, or FPGA overlay. Its architecture is
inspired by established devices, but its configuration format is its own.

\subsection{Host, guest and application}
The physical \textbf{host} is an XC7Z020, part \code{xc7z020clg400-1}. Vivado maps
bob's RTL onto host LUTs, registers, wires, RAM and DSP primitives. The
\textbf{guest} is bob: its configuration remains changeable after the host is
programmed. A \textbf{guest application} is a circuit mapped onto bob's resources.
\begin{lstlisting}[caption={Two separate implementation flows}]
HOST:  bob RTL -> Vivado synthesis/place/route -> AMD .bit
       -> board USB-JTAG -> physical XC7Z020 programmable logic
GUEST: application Verilog -> Yosys -> bob cells -> pack/place/route
       -> bob configuration -> Pico/PMODA JTAG -> running application
\end{lstlisting}
The Zynq ARM processing system, Linux, AXI and PYNQ notebooks are not required by
this flow. The board provides the programmable logic, clock, power and physical
I/O; the Pico supplies an independent guest configuration interface.

\subsection{Why two bitstreams and two JTAG interfaces?}
The board's built-in programming interface loads the host AMD bitstream. The
Pico talks to a second TAP implemented in bob's RTL on PMODA. These are separate
chains, not a pass-through connection. A guest chain cannot configure the host,
and an AMD bitstream cannot configure bob. A host rebuild changes the fabric
itself; a guest reload changes the circuit running on an unchanged fabric.

\subsection{What is original, reused and intentionally limited?}
The baseline TAP, boundary cells and CLB behaviour were inherited from the
hardware-tested predecessor. The project adds a central device schema, hardened
configuration plane, user clock control, hard-block wrappers, generated routing
and a synthesis-to-fabric flow. AMD UG470/473/474/479 guide configuration, RAM,
logic and DSP semantics. OpenFPGA/VPR guide routing-resource generation.
This does not establish complete compatibility with AMD device primitives,
bitstreams, BSDL descriptions, or every IEEE compliance requirement.
\src{REUSE.md}

\section{Physical platform and board connections}
Both Pico and PYNQ-Z2 use 3.3 V signal levels for this link. They require a common
ground. With each board powered separately, their power rails should not be
joined merely to make the JTAG connection. Verify the installed firmware's GPIO
assignment before wiring; the project documents the following mapping.
\begin{center}
\begin{tabular}{llll}
\toprule
Signal & Pico GPIO & PMODA pin & XC7Z020 package pin\\
\midrule
TCK & GP18 & 7 & U18\\
TMS & GP19 & 1 & Y18\\
TDI & GP16 & 2 & Y19\\
TDO & GP17 & 3 & Y16\\
Ground & GND & 5 or 11 & Ground\\
\bottomrule
\end{tabular}
\end{center}
TDI is driven by the probe into the target; TDO returns from target to probe.
These names are not crossed like UART TX/RX. U18 is a clock-capable input used
for TCK. The top instantiates separate BUFGs for TCK and the 125 MHz system clock
on H16; simulation bypasses these vendor primitives using \code{SIMULATION}.

\begin{center}
\begin{tabular}{llll}
\toprule
Board signal & Package pin & Logical pad & Edge\\
\midrule
SW0 & M20 & 6 & West\\
SW1 & M19 & 8 & West\\
BTN0 & D19 & 10 & West\\
BTN1 & D20 & 12 & West\\
BTN2 & L20 & 0 & South\\
BTN3 & L19 & 1 & South\\
LD0 & R14 & 7 & East\\
LD1 & P14 & 9 & East\\
LD2 & N16 & 11 & East\\
LD3 & M14 & Not a guest pad & DONE indicator\\
\bottomrule
\end{tabular}
\end{center}
Unused physical mappings have zero input and no external output connection;
their logical pads remain accessible through boundary scan. Raw switches and
buttons are asynchronous inputs. False-path constraints on these ports do not
synchronise or debounce them; sequential applications must respect that fact.
\src{hw/src/top/bob_top.v}
\src{hw/constr/pynq_z2.xdc}
\src{docs/wiring.md}

\section{Repository architecture and responsibility boundaries}
\begin{longtable}{p{0.23\textwidth}p{0.70\textwidth}}
\toprule
Area & Responsibility\\\midrule\endhead
\code{hw/} & Self-contained Windows/Vivado hand-off: ordered sources, configuration,
RTL, generated fabric, testbenches, constraints and Tcl scripts.\\
\code{tools/bob/} & Device description, VPR architecture and graph, generator,
cycle model, chain packing, synthesis, equivalence simulation and placement.\\
\code{host/} & DirtyJTAG USB transport, configuration operations, design API,
routing, board loading and milestone tests.\\
\code{sim/} & Vector generation, simulator drivers, lint and mutation testing.\\
\code{tests/} & Python tests for schema, layout, models, synthesis, reports and
Tcl build behaviour.\\
\code{examples/} & Six small Verilog applications used by the M8 flow.\\
\code{docs/} & Specifications, historical and current reports, wiring,
architecture visualisations, checklists and hardware evidence.\\
\code{release/} & Frozen hand-off bundles and restored historical tools.\\
\bottomrule
\end{longtable}
All directory names in this table are under
\nolinkurl{/Users/sk/work/bob/bob_full_v1/}.
\code{sources.f} is shared by Vivado, simulation and lint, preventing divergent
source lists. \code{device.py} owns architectural constants; generated JSON and
Verilog headers carry those constants to software and RTL consumers. Generated
files should not be fixed by hand: change the authority and regenerate instead.
\src{hw/sources.f}
\src{tools/bob/device.py}

\section{The active device and its hierarchy}
\label{sec:device}
\begin{center}
\begin{tabular}{ll}
\toprule
Property & Active M7 profile\\\midrule
Device name & \code{bob6x4}\\
VPR grid / interior & $8\times6$ / $6\times4$\\
CLBs & 16, one basic logic element (BLE) each\\
Memory & $2\times1024\times18$ bits\\
DSP & 2 cascaded DSP48E1-style slices\\
Logical pads / boundary cells & 20 / 40\\
Routing & L4 unidirectional, $W=24$, Wilton $F_s=3$\\
Programmable routing muxes & 1095\\
Configuration & 4216 bits, CRC-32C checked\\
IR width / IDCODE / USERCODE & 6 / \code{0xABEEF093} / \code{0x00000007}\\
\bottomrule
\end{tabular}
\end{center}
Coordinates increase eastward in $x$ and northward in $y$. Empty corners surround
a ring of I/O tiles. The four CLB columns are $x=1,2,4,5$. BRAM occupies column
$x=3$, and DSP occupies column $x=6$. Each hard block occupies two rows, rooted
at $y=1$ or $y=3$. Those footprint heights describe guest geometry, not the host
placement of inferred RAMB18 or DSP48E1 primitives.
\begin{lstlisting}[caption={Logical core; upper hard-block rows are continuations}]
       x=1  x=2  x=3    x=4  x=5  x=6
 y=4   CLB  CLB  BRAM1  CLB  CLB  DSP1
 y=3   CLB  CLB  BRAM1  CLB  CLB  DSP1
 y=2   CLB  CLB  BRAM0  CLB  CLB  DSP0
 y=1   CLB  CLB  BRAM0  CLB  CLB  DSP0
 Carry flows upward inside each CLB column.
 DSP0 PCOUT flows to DSP1 PCIN.
\end{lstlisting}
The frozen larger profile has a $10\times10$ VPR grid, an $8\times8$ interior,
48 CLBs, 32 pads, height-four hard blocks, and 9400 configuration bits. Its
IDCODE is \code{0x9BEEF093}. It is not interchangeable with the active profile.
The documented reason for downsizing was excessive synthesis runtime, not a
measured proof that the larger architecture could never fit.

\subsection{Top-down RTL composition}
\begin{lstlisting}[caption={Main hardware hierarchy}]
bob_top                       board pins and two BUFG clock buffers
  bob_fpga                    board-independent integration
    jtag_tap6                 TAP state machine and instruction decode
    cfg_ctrl                  CRC, count, status and startup sequencing
    cfg_tile_sr               configuration shift and shadow registers
    capture_chain             CLB-output observation
    clock_ctrl                sysclk-domain user enable and synchronizers
    bsc_cell x 40             boundary scan
    bram_jtag                 USER4 command engine
    dsp_jtag                  private DSP drive/readback register
    bob_fabric                generated routing and block instances
      bob_mux x 1095          configurable routing choices
      clb x 16                LUT, carry and flip-flop
      bram_block x 2          drive selection plus bram_core
      dsp_block x 2           drive selection plus dsp_core
\end{lstlisting}
\src{hw/src/fabric/bob_fpga.v}
\src{hw/src/generated/bob_fabric.v}

\section{Configurable logic block: Boolean logic to one stored bit}
\label{sec:clb}
\subsection{A lookup table is programmable truth-table memory}
For $K$ input bits, define the address
\[
a=\sum_{j=0}^{K-1}i_j2^j,\qquad O_6=\mathrm{INIT}[a].
\]
The input vector selects one of $2^K$ stored truth-table bits. With $K=6$, INIT
has 64 bits. Bob generalises the implementation to other supported $K$ values;
K=4 is also exercised in the verification flow. A LUT does not execute software:
its input-dependent selection is combinational hardware.

The secondary output is
\[
O_5=\mathrm{INIT}\left[\sum_{j=0}^{K-2}i_j2^j\right].
\]
For K=6 it uses the lower 32 INIT bits and ignores input 5. O5 and O6 share
storage; they are not two arbitrary independent six-input functions. Fracturing
must preserve that relationship when packing multiple uses into one CLB.

For a two-input AND with $i_0=A$, $i_1=B$, the low four truth-table bits are
$0,0,0,1$, hence hexadecimal 8. Replicating this nibble across all 64 bits gives
\code{64'h8888888888888888}, making the result independent of the other inputs.
For XOR the repeated nibble is 6. This address convention must agree in source
mapping, bit generation, the Python model and RTL.

\subsection{Carry arithmetic}
With propagate $P=O_6$ and generate input $D_I$ selected from $i_0$ or O5,
\[
C_{out}=\begin{cases}C_{in},&P=1\\D_I,&P=0,\end{cases}
\qquad S=P\oplus C_{in}.
\]
For addition set $P=A\oplus B$ and $D_I=A$. If the input bits differ, carry
propagates; if equal, they determine the new carry. The LUT supplies propagate,
the MUXCY-style expression supplies carry, and XORCY-style logic supplies sum.
Dedicated northward carry connections avoid routing every carry bit through
general-purpose channels. Crossing between columns requires explicit logic and
routing; it is not an invisible wraparound direct.

\subsection{Register priority and output selection}
Each CLB contains one register $q$. At a system-clock edge its update priority is:
\begin{enumerate}
\item If GSR is asserted, load the configured reset/initial value.
\item Otherwise, if either GWE or the user enable GCE is low, hold.
\item Otherwise, if enabled routed SR is asserted, load the reset value.
\item Otherwise, if effective CE is asserted, load the combinational datapath.
\item Otherwise hold.
\end{enumerate}
The main output selects $q$ or the combinational datapath. Carry mode selects
the sum; non-carry mode selects O6 or O5. Reset beats local CE, while GSR takes
priority independently of the user enable. Reset values zero and one provide
FDRE-like and FDSE-like behaviour respectively.
\src{hw/src/clb/lutk.sv}
\src{hw/src/clb/clb.sv}

\section{The active device and its hierarchy}
\label{sec:device}
\begin{center}
\begin{tabular}{ll}
\toprule
Property & Active M7 profile\\\midrule
Device name & \code{bob6x4}\\
VPR grid / interior & $8\times6$ / $6\times4$\\
CLBs & 16, one basic logic element each\\
Memory & $2\times1024\times18$ bits\\
DSP & 2 cascaded DSP48E1-style slices\\
Logical pads / boundary cells & 20 / 40\\
Routing & L4 unidirectional, $W=24$, Wilton $F_s=3$\\
Routing nodes / muxes / directs / undriven & 1564 / 1095 / 60 / 77\\
Configuration & 4216 bits, CRC-32C checked\\
IR width / IDCODE / USERCODE & 6 / \code{0xABEEF093} / \code{0x00000007}\\
\bottomrule
\end{tabular}
\end{center}
Coordinates increase eastward in $x$ and northward in $y$. Empty corners surround
an I/O ring. CLB columns are $x=1,2,4,5$; BRAM occupies column $x=3$ and DSP
column $x=6$, each hard block spanning two rows rooted at $y=1$ and $y=3$. These
heights are guest geometry; they do not describe host placement of inferred
RAMB18 or DSP48E1 primitives.
\begin{lstlisting}[caption={Logical core; upper hard-block rows are continuations}]
       x=1  x=2  x=3    x=4  x=5  x=6
 y=4   CLB  CLB  BRAM1  CLB  CLB  DSP1
 y=3   CLB  CLB  BRAM1  CLB  CLB  DSP1
 y=2   CLB  CLB  BRAM0  CLB  CLB  DSP0
 y=1   CLB  CLB  BRAM0  CLB  CLB  DSP0
 Carry flows north inside each CLB column; DSP0 PCOUT feeds DSP1 PCIN.
\end{lstlisting}
The frozen larger profile has a $10\times10$ VPR grid, an $8\times8$ interior,
48 CLBs, 32 pads, height-four hard blocks, a 9400-bit chain and IDCODE
\code{0x9BEEF093}. It is not interchangeable with the active profile. The
documented reason for downsizing was excessive Vivado synthesis runtime, not a
proof that the larger fabric could never fit.

\subsection{Top-down RTL composition}
\begin{lstlisting}[caption={Main hardware hierarchy}]
bob_top                       board pins and two BUFG clock buffers
  bob_fpga                    board-independent integration
    jtag_tap6                 TAP state machine and instruction decode
    cfg_ctrl                  CRC, count, status, startup sequencing
    cfg_tile_sr               configuration shift + shadow registers
    capture_chain             CLB-output observation register
    clock_ctrl                sysclk-domain user enable + synchronizers
    bsc_cell x 40             boundary scan cells
    bram_jtag / dsp_jtag      USER4 engine / private DSP register
    bob_fabric                generated routing and block instances
      bob_mux x 1095          configurable routing choices
      clb x 16                LUT, carry, flip-flop
      bram_block x 2          pin muxes plus bram_core
      dsp_block x 2           pin muxes plus dsp_core
\end{lstlisting}
\src{hw/src/fabric/bob_fpga.v}
\src{hw/src/generated/bob_fabric.v}

\section{The configurable logic block, bit by bit}
\label{sec:clb}
\subsection{A lookup table is programmable truth-table memory}
For $K$ input bits, with address $a=\sum_{j=0}^{K-1}i_j2^j$,
\[ O_6=\mathrm{INIT}[a], \qquad
   O_5=\mathrm{INIT}\Bigl[\textstyle\sum_{j=0}^{K-2}i_j2^j\Bigr]. \]
The active profile uses $K=6$, so INIT is 64 bits; the flow also exercises $K=4$.
A LUT executes no software: its inputs address stored truth-table bits
combinationally. The secondary output O5 shares storage with O6 (for $K=6$, the
low 32 INIT bits over $i_4..i_0$); it is not an independent function, and
packing must preserve that relationship.
For a two-input AND with $i_0=A$, $i_1=B$, the low four truth-table bits are
$0,0,0,1$ (hexadecimal 8); repeating that nibble yields
\code{64'h8888888888888888}, independent of the remaining inputs. XOR repeats 6.
This address convention must agree across mapping, bit generation, the Python
model and the RTL.

\subsection{Carry arithmetic}
With propagate $P=O_6$ and generate input $D_I$ selected from $i_0$ or O5,
\[ C_{out}=\begin{cases} C_{in}, & P=1\\ D_I, & P=0\end{cases}
   \qquad S=P\oplus C_{in}. \]
For addition, the LUT computes $P=A\oplus B$ and $D_I=A$: equal input bits
produce a defined carry, differing bits propagate it. Dedicated northward carry
connections keep long chains off general routing; crossing to another column
requires explicit routed logic, not an invisible wraparound.

\subsection{Register priority and output selection}
Each CLB has one register. At a user-enabled clock edge the priority is:
GSR (load INIT value); otherwise hold if GWE or GCE is low; otherwise routed SR
if enabled (load reset value, beating CE); otherwise load the datapath if CE is
effective; otherwise hold. The main output selects register or combinational
data; carry mode selects the sum, otherwise O6 or O5. Reset values 0 and 1 give
FDRE-like and FDSE-like behaviour.
\src{hw/src/clb/lutk.sv}
\src{hw/src/clb/clb.sv}

\section{Routing: from VPR graph to muxes and configuration bits}
\label{sec:routing}
Bob's interconnect is generated, not hand-drawn, following the OpenFPGA method:
\begin{enumerate}
\item \code{tools/bob/device.py} writes a VPR architecture derived from
OpenFPGA's \code{k6_frac_N10_tileable_adder_chain_dpram8K_dsp36_fracff_40nm.xml}:
a tileable layout with an I/O ring and empty corners, one L4 unidirectional
segment, Wilton switch pattern with $F_s=3$, connection fractions
$fc_{in}=0.15$, $fc_{out}=0.10$, and carry/clock pins excluded from general
routing.
\item \code{tools/bob/vpr_rrgraph.sh} runs VPR\,9 in the OpenFPGA Docker image,
routes a trivial benchmark, and writes the graph with a stamp holding the
architecture SHA-256, image digest and command.
\item \code{device.py} then refuses to run if the stamp does not match the
current architecture text and channel width, so RTL can never silently diverge
from the graph it claims to implement.
\end{enumerate}
Every CHANX/CHANY/IPIN node with more than one driver becomes a \code{bob_mux}
whose select comes from the chain:
\begin{center}
\begin{tabular}{ll}
\toprule
Mux output & Select encoding\\\midrule
CHANX / CHANY & 0 = constant 0; $1..N$ = driving inputs in node-id order\\
IPIN & 0 = constant 0; 1 = constant 1; $2..N{+}1$ = driving inputs\\
\bottomrule
\end{tabular}
\end{center}
An all-zero chain therefore yields a loop-free, dark fabric. An IPIN whose only
driver is a single OPIN is not a mux at all but a wire; all 60 directs are of
this kind (the CLB carry chain and the DSP cascade). Undriven nodes read zero.

\subsection{Where the routing bits live}
The chain order is: control tile, then grid tiles row-major from $(0,0)$, then
padding to a byte boundary. A grid tile stores its block fields first, then its
routing-mux selects by ascending rr-node id. Examples from
\code{tools/bob/device.json}:
\begin{center}
\begin{tabular}{llll}
\toprule
Tile & Chain bits & Width & Contents\\\midrule
\code{ctrl} & 0--7 & 8 & \code{clk_mode}, \code{clk_div}, reserved\\
\code{t_x1y1} (CLB) & 192--352 & 161 & 71 block bits + 90 routing bits\\
\code{t_x3y1} (BRAM0) & 519--654 & 136 & 8 block bits + 128 pin selects\\
\code{t_x6y1} (DSP0) & 984--1203 & 220 & 16 block bits + 204 selects\\
\code{tail} & 4211--4215 & 5 & reserved, write 0\\
\bottomrule
\end{tabular}
\end{center}
The 71 CLB block bits are the 64-bit INIT plus seven flags (\code{ff_en},
\code{ff_rstval}, \code{ff_ce_en}, \code{ff_sr_en}, \code{cy_en},
\code{cy_di_sel}, \code{ff_d_sel}). Tile widths vary with the muxes each tile
owns, which is why nominally identical CLB tiles differ slightly.
A documented pitfall shaped the pin plan: VPR's \code{pinlocations
pattern="spread"} places pins on the top and bottom faces \emph{inside} tall
blocks, where \code{through_channel="false"} leaves no channel, producing
unroutable pins. Hard-block pins are therefore placed round-robin on the left
and right of every row, with clock, carry and cascade pins on the block ends.
\src{tools/bob/device.py}
\src{tools/bob/fabric_gen.py}

\section{Routing: from VPR's graph to muxes and chain bits}
\label{sec:routing}
Bob's interconnect is generated, not hand-drawn, following the OpenFPGA method:
\begin{enumerate}
\item \code{tools/bob/device.py} writes a VPR architecture derived from
OpenFPGA's \code{k6_frac_N10_tileable_adder_chain_dpram8K_dsp36_fracff_40nm.xml}:
tileable layout, I/O ring with empty corners, one L4 unidirectional segment,
Wilton switch pattern with $F_s=3$, connection fractions
$fc_{\mathrm{in}}=0.15$ and $fc_{\mathrm{out}}=0.10$, and carry/clock pins
excluded from general routing.
\item \code{tools/bob/vpr_rrgraph.sh} runs VPR~9 in the OpenFPGA Docker image,
routes a trivial benchmark, and writes the routing-resource graph together with
a stamp holding the architecture SHA-256, image digest and command.
\item Regeneration then fails if the stamp does not match the current
architecture text or channel width, so RTL can never silently diverge from the
graph it claims to implement.
\end{enumerate}
Every CHANX/CHANY/IPIN node with more than one driver becomes a \code{bob_mux}
whose select comes from the configuration chain:
\begin{center}
\begin{tabular}{ll}
\toprule
Mux output & Select encoding\\\midrule
CHANX / CHANY & 0 = constant 0; $1..N$ = driving inputs in node-id order\\
IPIN & 0 = constant 0; 1 = constant 1; $2..N{+}1$ = driving inputs\\
\bottomrule
\end{tabular}
\end{center}
An all-zero chain is therefore a dark, loop-free fabric. An input pin whose only
driver is a single output pin is not a mux at all but a wire; all 60 directs
are of this kind (the carry chain and the DSP cascade). Undriven nodes read
constant zero.

\subsection{Where the routing bits live in the chain}
The chain order is: control tile, grid tiles row-major from $(0,0)$, then
padding to a byte boundary. A grid tile stores its block fields first, then its
routing-mux selects by ascending rr-node id. Examples from
\code{tools/bob/device.json}:
\begin{center}
\begin{tabular}{llll}
\toprule
Tile & Chain bits & Width & Contents\\\midrule
\code{ctrl} & 0--7 & 8 & \code{clk_mode}, \code{clk_div}, reserved\\
\code{t_x1y1} (CLB) & 192--352 & 161 & 71 block bits + 90 routing bits\\
\code{t_x3y1} (BRAM0) & 519--654 & 136 & 8 block bits + 128 pin selects\\
\code{t_x6y1} (DSP0) & 984--1203 & 220 & 16 block bits + 204 selects\\
\code{tail} & 4211--4215 & 5 & reserved, write 0\\
\bottomrule
\end{tabular}
\end{center}
The 71 CLB block bits are the 64-bit INIT plus seven flags: \code{ff_en},
\code{ff_rstval}, \code{ff_ce_en}, \code{ff_sr_en}, \code{cy_en},
\code{cy_di_sel}, \code{ff_d_sel}. Tile widths vary with the muxes each tile
owns, which is why nominally identical CLB tiles differ slightly in width.
A documented pitfall shaped the pin plan: VPR's \code{pinlocations
pattern="spread"} places pins on the top and bottom faces \emph{inside} tall
blocks, where \code{through_channel="false"} leaves no channel, producing
unroutable pins. Hard-block pins are therefore placed round-robin on the left
and right of every row, with clock, carry and cascade pins on the block ends.
\src{tools/bob/device.py}
\src{tools/bob/fabric_gen.py}

\section{Hard blocks: block memory and multiply--accumulate}
\label{sec:hard}
\subsection{BRAM: UG473 behaviour on an inferred RAMB18}
\code{bram_core.v} implements a $1024\times18$ true-dual-port memory following
AMD UG473 semantics, trimmed to one width mode. Per-port configuration bits
cover write mode (WRITE\_FIRST, READ\_FIRST, NO\_CHANGE), optional output
register (DOA/DOB\_REG style), and per-port enable and reset. The memory array
is written in Vivado's true-dual-port READ\_FIRST template so that synthesis
infers a RAMB18 primitive rather than distributed LUT RAM. WRITE\_FIRST and
NO\_CHANGE are derived from the READ\_FIRST template by registers placed after
the RAM output; this is what makes the write mode a configuration bit rather
than a synthesis attribute.
GSR clears output latches and registers to zero but never the memory contents.
Contents load through port~A from the USER4 engine while GWE is low, consistent
with UG470's separation of configuration from block-memory contents.
Eight configuration bits per block select write modes, output registers, and
whether each port's pins come from the fabric or from the 64-bit USER4 drive
word. Hard-block pin selects are routing bits: each pin is a node in the rr
graph, so a BRAM or DSP pin can be driven from fabric tracks, constants, or the
JTAG drive word.
\begin{center}
\begin{tabular}{lll}
\toprule
USER4 command & Code & Behaviour\\\midrule
LOAD\_PTR & 1 & address pointer $\leftarrow$ payload\\
WRITE & 2 & write word at pointer, increment (GWE = 0 only)\\
READ & 3 & read word at pointer, increment (GWE = 0 only)\\
SET\_DRIVE & 4 & selected block's drive word $\leftarrow$ payload\\
SELECT & 5 & target block index $\leftarrow$ payload\\
\bottomrule
\end{tabular}
\end{center}
\subsection{DSP: UG479 behaviour, trimmed}
\code{dsp_core.v} implements one DSP48E1-style slice with static configuration
instead of dynamic OPMODE/INMODE/ALUMODE pins. Inputs A (25), B (18), C (48) and
D (25) are signed; a pre-adder computes $AD = D \pm A$; the multiplier computes
$M = AD\times B$ over 43 bits. The output follows one of four modes:
$P=M$; $P=M+C$; $P=P+M$ (accumulator, always fed back from the P register, so
no combinational loop exists when PREG is zero); and cascade accumulation,
$P=(\code{PCIN} >>> 17)+M$, using the 17-bit shift familiar from UG479.
Register stages AREG, BREG, CREG, DREG, MREG and PREG are configuration bits;
CE and reset are grouped in four pairs, with C sharing P's controls, and reset
beating enable as UG479's synchronous resets do. GCE, GWE and GSR act on every
register exactly as on the CLB. Written in UG479 inference style, the slice maps
to a host DSP48E1 primitive; two slices exist, with slice~0's PCOUT feeding
slice~1's PCIN through a direct.
\src{hw/src/tiles/bram_core.v}
\src{hw/src/tiles/dsp_core.v}

\section{Configuration: the scan chain, CRC and startup, bit by bit}
\label{sec:config}
\subsection{JTAG instructions}
The TAP implements a six-bit instruction register using AMD 7-series codes
compiled from OpenOCD, openFPGALoader and the XC7Z020 BSDL. Capture-IR returns
\code{\{DONE, INIT_B, COMMITTED, CRC_ERR, 0, 1\}}, so a plain IR scan reads
configuration status; the two least significant bits are the IEEE 1149.1
requirement. INTEST (\code{000111}) and the DSP instruction (\code{101000}) are
private codes not found in AMD's numbering.
\begin{center}
\begin{tabular}{lll}
\toprule
Code & Instruction & Data register\\\midrule
\code{000001} & SAMPLE & boundary ring, observe pins\\
\code{100110} & EXTEST & boundary ring, drive pads\\
\code{000111} & INTEST & boundary ring, drive fabric inputs\\
\code{000010} & USER1 & 32-bit control (ce, sr, cin, step, autostep)\\
\code{000011} & USER2 / CFG\_CTRL & 64-bit status and expected CRC\\
\code{100010} & USER3 / CAPTURE & 16 CLB outputs, read-only\\
\code{000100} & CFG\_OUT & chain, readback, never commits\\
\code{000101} & CFG\_IN & chain, commits when CRC and length pass\\
\code{001000} & USERCODE & 32-bit milestone code\\
\code{001001} & IDCODE & 32-bit identification\\
\code{001011} & JPROGRAM & clears configuration, re-asserts GSR/GTS\\
\code{001100} & JSTART & steps startup in Run-Test/Idle\\
\code{100011} & USER4 & BRAM contents, drive and SELECT\\
\code{101000} & DSP & private DSP drive and P readback\\
\code{111111} & BYPASS & one bit\\
\bottomrule
\end{tabular}
\end{center}
\subsection{Commit rules, CRC and the startup sequence}
Every CFG\_IN scan computes a serial CRC-32C (Castagnoli polynomial
\code{0x1EDC6F41}, reflected \code{0x82F63B78}, init and final XOR
\code{0xFFFFFFFF}) over the shifted bits, LSB first, and counts them. At
Update-DR the chain commits only if the count equals 4216 and the complement of
the running CRC equals the expected value written beforehand through CFG\_CTRL
behind the key byte \code{0xC5}. A rejected scan raises \code{CRC\_ERR} or
\code{LEN\_ERR} and leaves the previous configuration intact. CFG\_OUT shifts
the shadow chain out and never commits.
The status word exposed through CFG\_CTRL is \code{\{version, done, gwe, gts,
gsr, committed, len_err, crc_err, crc_ok, count[15:0], crc[31:0]\}}.
Startup follows UG470's order: JPROGRAM re-asserts GSR and GTS and clears
committed and DONE; each JSTART clock in Run-Test/Idle advances one phase,
releasing GSR, then GTS, then GWE, then DONE. The chain is LSB-first: the first
bit shifted in on TDI is chain bit 0, and the tail's reserved bits are last.
The commit gate currently does not check GWE: a correct live reload commits
while the fabric runs. That is deliberate for fast design swaps and is a known
divergence from UG470's full-configuration model; it is also the reason a
blanket false-path from configuration flops would describe an operating mode
the RTL still permits.
\begin{lstlisting}[caption={Host load sequence (host/cfgplane.py)}]
1. IR JPROGRAM          clear chain, GSR/GTS asserted, DONE low
2. CFG_CTRL write       (0xC5 << 56) | expected CRC-32C
3. CFG_IN shift         4216 bits, bulk transfer allowed
4. CFG_CTRL read        require committed, crc_ok, count = 4216
5. CFG_OUT read         must equal the sent chain
6. JSTART + 12 TCK      startup sequence
7. IR capture           require DONE = 1
\end{lstlisting}
A host-side \code{.bobc} container (magic \code{BOBC}, version, device name,
width, CRC, payload) is validated in full before any hardware is touched.
\src{hw/src/core/cfg_ctrl.v}
\src{docs/bitstream-format.md}
\src{tools/bob/chainbits.py}

\section{User clock, CDC and boundary scan}
\label{sec:clock}
Every fabric flop runs on the 125 MHz board clock through a BUFG and may change
only on an edge where GCE is high. Clock control realises the user clock as an
enable rather than a derived clock, following UG949's recommendation: JTAG mode
produces one pulse per TCK rising edge while USER1 ce is set (plus manual step
or INTEST autostep pulses), and free-running mode produces one pulse every
$2^{\mathrm{clk\_div}+8}$ sysclk cycles, at most 488 kHz. All TCK-domain inputs
pass through two-flop ASYNC\_REG synchronizers. The XDC therefore grants
120-cycle multicycles on fabric flop-to-flop paths, justified by the free-running
minimum of 256 cycles and by a TCK ceiling of about 1\,MHz.
These are operating contracts, not self-enforcing properties. The review
identified three gaps that remain open in the source: nothing enforces one
enable source at a time (a manual step on a TCK rising edge and an autostep on
the falling edge can produce two GCE pulses roughly half a TCK period apart,
about 62 sysclk cycles at 1 MHz); a mode-switch commit can place two pulses
close together; and \code{dirtyjtag.py} applies any requested frequency without
a ceiling check. Consequences depend on guest designs; they are not known
hardware failures on this build.
Boundary scan uses 40 cells, two per pad: the first NPAD cells observe the
fabric's pad outputs (GTS forces them to zero before startup), the next NPAD
cells drive fabric inputs. SAMPLE is transparent observation; EXTEST drives
pads; INTEST drives fabric inputs, with autostep providing one user-clock edge
per scan. Test-Logic-Reset is consumed synchronously only, after a documented
combinational glitch caused real hardware failures in the predecessor.
\src{hw/src/core/clock_ctrl.v}
\src{host/dirtyjtag.py}

\section{Guest design flow: from Verilog to a running fabric}
\label{sec:flow}
\begin{lstlisting}[caption={M8 flow: each step has a proof obligation}]
tools/bob/synth.py examples/counter.v
  Yosys synth_xilinx-style flow restricted to bob cells:
  $lut, BOB_FDRE/FDSE, BOB_ADD, BOB_BRAM18, BOB_DSP.
  Fails if a foreign cell remains or the design uses >1 clock.
tools/bob/equiv.py examples/counter.v
  Source vs synthesised netlist in one iverilog bench, 300 random
  cycles, outputs compared before and after every edge; the source
  trace is saved for later reuse.
tools/bob/place.py counter --check
  Deterministic packing/placement, then rr-graph routing; --check
  replays the source trace through the cycle model.
sim/run_synth_sim.sh
  The same placed bitstreams on the complete fabric RTL (tb_synth).
\end{lstlisting}
Packing merges an adder and its flip-flop into one carry-mode CLB, merges a LUT
feeding only one FF, keeps extra loads of small LUTs on O5, and runs carry
chains up a CLB column from a generator row, so the global carry-in never enters
a design; chains longer than a column continue through a tap CLB in the next
carry column. Placement is deliberately simple and deterministic; VPR takes over
packing, placement and routing at M9. Known M8 limits recorded in the plan: no
asynchronous resets or latches, one clock, no mid-chain carry tap, no BRAM
output-register or DSP register absorption.
\subsection{Routing a guest net}
\code{host/bitstream.py} routes nets by breadth-first search over the rr graph,
sharing tracks among fanout and reordering rip-ups when a net fails;
\code{tools/bob/model.py} evaluates exactly the muxes and directs the RTL has,
so model and RTL agree by construction. The equivalence harness exists because
mapping bugs are silent: a wrong INIT addressing convention or mis-ordered
carry bit still simulates, just wrongly.
\src{tools/bob/synth.py}
\src{tools/bob/place.py}
\src{host/bitstream.py}

\section{Verification: layers, mutants and evidence}
\label{sec:verify}
\begin{center}
\begin{tabular}{lll}
\toprule
Layer & Mechanism & Recorded result\\\midrule
Config plane & tb\_cfg vectors, per-guard mutants & 180 checks; 5/5 killed\\
CLB sweep & 128 flag combinations vs model & 7040 checks (K=6)\\
BRAM & mode/register sweep vs model & 5292 checks\\
DSP & opmode/register sweep + cascade & 1344 checks\\
Fabric (48-CLB) & tb\_bob incl.\ random routed netlists & 852 checks\\
K=4 substitution & full suite at K=4 & 722 checks\\
Synthesis & 6 examples, 3-way equivalence & tb\_synth 486\\
Fabric mutants & ce/sr/divider/gce mutants & 25/25 killed\\
pytest & device, layout, Tcl, reports, synth & 88 passed\\
\bottomrule
\end{tabular}
\end{center}
These are recorded results from the plan and logs; not rerun for this report.
The independent software model is the oracle: RTL checks compare against it,
and it is itself exercised by mutation tests that prove each guard's absence
would be caught. Milestone hardware tests always run the full earlier
regression first, so a new milestone cannot pass while breaking an old one.

\section{Host build flow and the Windows hand-off}
\label{sec:build}
Vivado runs on a separate Windows machine; \code{hw/} is the complete bundle.
\code{build.tcl} reuses an existing Vivado project, syncs it to
\code{sources.f}, fingerprints sources so it rebuilds only on real changes, and
errors out if implementation claims success without a bitstream. It writes
timing, utilisation and DRC reports plus a build-info fingerprint (tag, top,
IDCODE, USERCODE, part, source list, WNS) into \code{out/<tag>/}, which is
copied back to \code{docs/reports/<tag>/}. The fabric's combinational loops are
by construction, so LUTLP-1 is downgraded to a warning for fabric tops only,
via a pre-implementation and pre-bitstream Tcl hook. Synthesis uses
\code{RuntimeOptimized} because timing-driven synthesis of the looped fabric is
slow; this changes optimisation effort, not the meaning of subsequent static
timing analysis. On negative WNS the script warns, explains that static timing
is inconclusive on this design by construction, and still writes the bitstream.
\src{hw/scripts/build.tcl}
\src{hw/scripts/drc_waiver.tcl}
\src{tests/test_build_tcl.py}

\section{Hardware status and the evidence trail}
\label{sec:status}
\begin{longtable}{p{0.14\textwidth}p{0.46\textwidth}p{0.28\textwidth}}
\toprule
Milestone & Content & Evidence status\\\midrule\endhead
M0--M2 & Layout, schema, config plane & Vivado reports archived; board checks logged\\
M3 & 4$\times$4 fabric on the config plane & Report archived; board checks logged\\
M4 & User clock, routed CE/SR, K parameter & Report archived; WNS $+0.843$\,ns\\
M5 & BRAM tile & Board checks logged; utilisation report not copied back\\
M6 & DSP tile & Board regression logged passing; report missing\\
M7 & Generated fabric, 16-CLB profile & Simulated; hardware test not passed\\
M8 & Synthesis onto bob cells & Simulated; board test rides the M7 bitstream\\
M9--M13 & VPR PnR, bitgen, bring-up, frames & Planned\\\bottomrule
\end{longtable}
The last M7 log entries are three identification failures returning
\code{0xFFFFFFFF}, which per the project's own debugging table means nothing
answered on TDO: the board was not programmed, or the probe was disconnected.
They do not test the fabric, so they neither validate nor refute it; the most
recent passing board regression remains M6. Note that IDCODE/USERCODE identity
alone does not prove which rebuild is loaded, since rebuilds of the same
profile share both codes; a source fingerprint or bitstream hash would close
that gap in future acceptance checks.

\section{Open issues and recommended next steps}
\label{sec:issues}
\begin{enumerate}
\item \textbf{Version control.} The tree has no Git history; the M6 host tools
had to be rebuilt from a session transcript after an overwrite. Recommended:
initialise Git, tag milestones, and derive \code{release/} bundles from tags
(archive or worktree) rather than manual copies, keeping the Windows hand-off
workflow.
\item \textbf{M7 acceptance.} Obtain the 16-CLB Vivado reports, analyse any
negative WNS by actual startpoint and endpoint, then run the hardware
regression. Nothing should be waived before the real paths are inspected.
\item \textbf{Operating contracts as code.} Commit only while GWE = 0 (with
matching status and host interpretation), one enable source at a time, a TCK
ceiling enforced in \code{dirtyjtag.py}, consistency tests tying the multicycle
value to the divider and probe ceiling, and a \code{report_cdc} review of the
BRAM command handshake and quasi-static drive words.
\item \textbf{Report acceptance.} \code{test_reports.py} checks only folders
that already exist and asserts no slack sign. A separate acceptance target
should require the current build's reports, IDENTITY match, nonnegative slack
or a justified exception, expected RAMB18/DSP counts, and a logged hardware
pass; everyday checks stay usable before reports exist.
\item \textbf{Clearer dead-probe diagnosis.} \code{0xFFFFFFFF} currently
produces the same "wrong bitstream" message as a genuine mismatch; it should
say that nothing is driving TDO.
\item \textbf{Scaling.} Measure LUTRAM configuration memory and out-of-context
or incremental builds before growing the grid.
\end{enumerate}
\src{tests/test_reports.py}
\src{host/hwtest.py}

\section{How to reproduce the work}
\label{sec:repro}
On the Mac, \code{make check} regenerates device files, runs every simulation,
lint and pytest, and must be green before any hand-off; \code{make rrgraph}
rebuilds the routing graph in Docker after an architecture change; \code{make
hw} refreshes generated vectors inside \code{hw/}. On the Windows machine,
\code{build.tcl} rebuilds only what changed and can program the board. Back on
the Mac, \code{make hwtest M=Mx} runs the regression plus the milestone's
checks and appends \code{docs/hwtest/results.log}. Older bitstreams are tested
with the frozen tools under \code{release/}. This report compiles with:
\begin{lstlisting}[language={}]
cd /Users/sk/work/bob/bob_full_v1/docs
pdflatex bob_full_v1_report.tex   # run twice for the table of contents
\end{lstlisting}
No compilation was performed for this delivery.

\end{document}
